% GENERATED FILE. Edit canonical lessons, then run npm run generate:books.
% canonical-source-hash: fnv1a64:68979168384b17de
% canonical-lessons: JA-C268-ochitsuku, JA-C268-kyuujo, JA-C268-ikinokoru, JA-C268-tasukeau, JA-C268-anzen

\chapter{Keeping Calm}
\label{ch:ja-keeping-calm}
\hlchaptermodality{268}

\begin{chapteropening}
\textbf{By the end of this chapter:} \emph{I can say what to do in an emergency and who survived.}

Complete the last lesson of chapter 268: I can say what to do in an emergency and who survived.
\end{chapteropening}

\section[\texorpdfstring{\ja{おちつく} (ochitsuku)}{おちつく (ochitsuku)}]{\ja{おちつく} (ochitsuku) — to calm down; to keep calm}

\label{lesson:JA-C268-ochitsuku}



\begin{quote}
Before the new one: say the Japanese for a cry, then the Japanese for with relief.
\end{quote}

\subsection*{You'll want to know: \ja{おちつく}}
\textbf{\ja{おちつく}} — \emph{ochitsuku} — \textquotedblleft{}to calm down; to keep calm\textquotedblright{}.

\begin{quote}
\textbf{\ja{じこが} \ja{あったら、おちついて} \ja{ください。}} — \emph{jiko ga attara, ochitsuite kudasai} — \textquotedblleft{}If there is an accident, please keep calm.\textquotedblright{}
\end{quote}

Polite \textbf{\ja{おちつきます}}, \emph{ochitsukimasu}; a godan verb. Past (polite): \textbf{\ja{ははが} \ja{おちつきました。}}, \emph{haha ga ochitsukimashita}.

\begin{cousinweb}
\textbf{\ja{おちつく}} (\emph{ochitsuku}) is a native Japanese word (wago), inherited rather than borrowed from Chinese, built from \emph{ochi}, \textquotedblleft{}falling\textquotedblright{} (from \emph{ochiru}), and \emph{tsuku}, \textquotedblleft{}to settle\textquotedblright{} (JMdict, KANJIDIC2).
\end{cousinweb}

\begin{grammarlens}[title={What you've built}]
One more word for telling what happened.
\end{grammarlens}

\subsection*{Your turn}
\begin{itemize}
\raggedright
  \item \emph{Say it:} \emph{ochitsuku}
  \item \emph{Say it:} \emph{ochitsuku}, once more
  \item \emph{Say it:} \emph{jiko ga attara, ochitsuite kudasai} — If there is an accident, please keep calm.
  \item \emph{Recall:} say the Japanese for a cry, then the Japanese for with relief, then say \emph{ochitsuku} again
\end{itemize}

\begin{tcolorbox}[breakable,colback=teal!4,colframe=teal!35!black,title={Before you move on}]
What does \ja{おちつく} mean? (\textquotedblleft{}To calm down; to keep calm\textquotedblright{}.) Say it once more.
\end{tcolorbox}
\section[\texorpdfstring{\ja{きゅうじょ} (kyūjo)}{きゅうじょ (kyūjo)}]{\ja{きゅうじょ} (kyūjo) — a rescue; relief}

\label{lesson:JA-C268-kyuujo}



\begin{quote}
Before the new one: say the Japanese for with relief, then the Japanese for to calm down.
\end{quote}

\subsection*{You'll want to know: \ja{きゅうじょ}}
\textbf{\ja{きゅうじょ}} — \emph{kyūjo} — \textquotedblleft{}a rescue; relief\textquotedblright{}.

\begin{quote}
\textbf{\ja{みんなで} \ja{きゅうじょを} \ja{まっていました。}} — \emph{minna de kyūjo o matte imashita} — \textquotedblleft{}We all waited to be rescued.\textquotedblright{}
\end{quote}

\begin{cousinweb}
\textbf{\ja{きゅうじょ}} (\emph{kyūjo}) is Sino-Japanese (kango): two Chinese characters read with their on'yomi, readings borrowed from Chinese. They are \emph{kyū}, \textquotedblleft{}rescue\textquotedblright{} (Middle Chinese \emph{kjuwH}) and \emph{jo}, \textquotedblleft{}help\textquotedblright{} (\emph{dzrjoH}) (KANJIDIC2; Middle Chinese in Baxter's transcription).
\end{cousinweb}

\begin{grammarlens}[title={What you've built}]
One more word for telling what happened.
\end{grammarlens}

\subsection*{Your turn}
\begin{itemize}
\raggedright
  \item \emph{Say it:} \emph{kyūjo}
  \item \emph{Say it:} \emph{kyūjo}, once more
  \item \emph{Say it:} \emph{minna de kyūjo o matte imashita} — We all waited to be rescued.
  \item \emph{Recall:} say the Japanese for with relief, then the Japanese for to calm down, then say \emph{kyūjo} again
\end{itemize}

\begin{tcolorbox}[breakable,colback=teal!4,colframe=teal!35!black,title={Before you move on}]
What does \ja{きゅうじょ} mean? (\textquotedblleft{}A rescue; relief\textquotedblright{}.) Say it once more.
\end{tcolorbox}
\section[\texorpdfstring{\ja{いきのこる} (ikinokoru)}{いきのこる (ikinokoru)}]{\ja{いきのこる} (ikinokoru) — to survive; to come through}

\label{lesson:JA-C268-ikinokoru}



\begin{quote}
Before the new one: say the Japanese for to calm down, then the Japanese for a rescue.
\end{quote}

\subsection*{You'll want to know: \ja{いきのこる}}
\textbf{\ja{いきのこる}} — \emph{ikinokoru} — \textquotedblleft{}to survive; to come through\textquotedblright{}.

\begin{quote}
\textbf{\ja{じこで} \ja{ふたりが} \ja{いきのこりました。}} — \emph{jiko de futari ga ikinokorimashita} — \textquotedblleft{}Two people came through the accident alive.\textquotedblright{}
\end{quote}

Polite \textbf{\ja{いきのこります}}, \emph{ikinokorimasu}; a godan verb. Past (polite): \textbf{\ja{いぬが} \ja{いきのこりました。}}, \emph{inu ga ikinokorimashita}.

\begin{cousinweb}
\textbf{\ja{いきのこる}} (\emph{ikinokoru}) is a native Japanese word (wago), inherited rather than borrowed from Chinese, built from \emph{iki}, \textquotedblleft{}living\textquotedblright{} (from \emph{ikiru}), and \emph{nokoru}, \textquotedblleft{}to remain\textquotedblright{} (JMdict, KANJIDIC2).
\end{cousinweb}

\begin{grammarlens}[title={What you've built}]
One more word for telling what happened.
\end{grammarlens}

\subsection*{Your turn}
\begin{itemize}
\raggedright
  \item \emph{Say it:} \emph{ikinokoru}
  \item \emph{Say it:} \emph{ikinokoru}, once more
  \item \emph{Say it:} \emph{jiko de futari ga ikinokorimashita} — Two people came through the accident alive.
  \item \emph{Recall:} say the Japanese for to calm down, then the Japanese for a rescue, then say \emph{ikinokoru} again
\end{itemize}

\begin{tcolorbox}[breakable,colback=teal!4,colframe=teal!35!black,title={Before you move on}]
What does \ja{いきのこる} mean? (\textquotedblleft{}To survive; to come through\textquotedblright{}.) Say it once more.
\end{tcolorbox}
\section[\texorpdfstring{\ja{たすけあう} (tasukeau)}{たすけあう (tasukeau)}]{\ja{たすけあう} (tasukeau) — to help each other}

\label{lesson:JA-C268-tasukeau}



\begin{quote}
Before the new one: say the Japanese for a rescue, then the Japanese for to survive.
\end{quote}

\subsection*{You'll want to know: \ja{たすけあう}}
\textbf{\ja{たすけあう}} — \emph{tasukeau} — \textquotedblleft{}to help each other\textquotedblright{}.

\begin{quote}
\textbf{\ja{こまった} \ja{とき、みんなで} \ja{たすけあいます。}} — \emph{komatta toki, minna de tasukeaimasu} — \textquotedblleft{}In hard times we all help one another.\textquotedblright{}
\end{quote}

Polite \textbf{\ja{たすけあいます}}, \emph{tasukeaimasu}; a godan verb. Past (polite): \textbf{\ja{みんなが} \ja{たすけあいました。}}, \emph{minna ga tasukeaimashita}.

\begin{cousinweb}
\textbf{\ja{たすけあう}} (\emph{tasukeau}) is a native Japanese word (wago), inherited rather than borrowed from Chinese, built from \emph{tasuke}, \textquotedblleft{}helping\textquotedblright{} (from \emph{tasukeru}), and \emph{au}, \textquotedblleft{}to do together\textquotedblright{} (JMdict, KANJIDIC2).
\end{cousinweb}

\begin{grammarlens}[title={What you've built}]
One more word for telling what happened.
\end{grammarlens}

\subsection*{Your turn}
\begin{itemize}
\raggedright
  \item \emph{Say it:} \emph{tasukeau}
  \item \emph{Say it:} \emph{tasukeau}, once more
  \item \emph{Say it:} \emph{komatta toki, minna de tasukeaimasu} — In hard times we all help one another.
  \item \emph{Recall:} say the Japanese for a rescue, then the Japanese for to survive, then say \emph{tasukeau} again
\end{itemize}

\begin{tcolorbox}[breakable,colback=teal!4,colframe=teal!35!black,title={Before you move on}]
What does \ja{たすけあう} mean? (\textquotedblleft{}To help each other\textquotedblright{}.) Say it once more.
\end{tcolorbox}
\section[\texorpdfstring{\ja{あんぜん} (anzen)}{あんぜん (anzen)}]{\ja{あんぜん} (anzen) — safe; safety}

\label{lesson:JA-C268-anzen}



\begin{quote}
Before the new one: say the Japanese for to survive, then the Japanese for to help each other.
\end{quote}

\subsection*{You'll want to know: \ja{あんぜん}}
\textbf{\ja{あんぜん}} — \emph{anzen} — \textquotedblleft{}safe; safety\textquotedblright{}.

\begin{quote}
\textbf{\ja{うちの} \ja{なかは} \ja{あんぜんでした。}} — \emph{uchi no naka wa anzen deshita} — \textquotedblleft{}It was safe inside the house.\textquotedblright{}
\end{quote}

A na-adjective: \ja{な} before a noun.

\begin{cousinweb}
\textbf{\ja{あんぜん}} (\emph{anzen}) is Sino-Japanese (kango): two Chinese characters read with their on'yomi, readings borrowed from Chinese. They are \emph{an}, \textquotedblleft{}quiet\textquotedblright{} (Middle Chinese \emph{'an}) and \emph{zen}, \textquotedblleft{}whole\textquotedblright{} (\emph{dzjwen}) (KANJIDIC2; Middle Chinese in Baxter's transcription).
\end{cousinweb}

\begin{grammarlens}[title={What you've built}]
One more word for telling what happened.
\end{grammarlens}

\subsection*{Your turn}
\begin{itemize}
\raggedright
  \item \emph{Say it:} \emph{anzen}
  \item \emph{Say it:} \emph{anzen}, once more
  \item \emph{Say it:} \emph{uchi no naka wa anzen deshita} — It was safe inside the house.
  \item \emph{Recall:} say the Japanese for to survive, then the Japanese for to help each other, then say \emph{anzen} again
\end{itemize}

\begin{tcolorbox}[breakable,colback=teal!4,colframe=teal!35!black,title={Before you move on}]
What does \ja{あんぜん} mean? (\textquotedblleft{}Safe; safety\textquotedblright{}.) Say it once more.
\end{tcolorbox}
